\documentclass[11pt]{article}
\usepackage{preamble}

\begin{document}

\maketitle
\thispagestyle{fancy}

\section*{BLOQUE 1--COMBUSTIBLES}

\begin{multicols*}{2}

    \begin{equation}
        PCI = PCS + g_v \cdot H_{vap.agua}
    \end{equation}

    De donde $H_{vap.agua} = 540 \ \mathrm{Kcal/Kg}$

    \begin{equation}
        PCI = PCS - 540 \cdot H \ \text{(en Kcal/Kg)}
    \end{equation}

    De donde $H$ es el porcentaje de hidrógeno en el combustible (en tanto por ciento).

    \begin{equation}
        P_e = \frac{P}{V} = \frac{M \cdot g}{V} = \rho \cdot g
    \end{equation}

    \paragraph*{Viscosidad dinámica}

    \begin{equation*}
        1 Po = \frac{dina \cdot cm}{cm^{2} \cdot \frac{cm}{s}} = \frac{g}{cm \cdot s}
    \end{equation*}

    1 cp = 0,01 Po = $10^{-3}$ kg/m$\cdot$s

    \paragraph*{Viscosidad cinemática}

    \begin{equation*}
        1 St = \frac{dina \cdot cm}{g} = \frac{cm^2}{s}
    \end{equation*}
    1 cSt = 0,01 St = $10^{-6}$ m$^2$/s

    \paragraph*{Índice de Wobbe}
    \begin{equation}
        W = \frac{PC}{\sqrt{d}}
    \end{equation}

    \paragraph*{Índice de Delbourg}
    \begin{equation}
        C = U \cdot \frac{H_2 + 0.7CO + 0.3CH_4 - V\sum_{i=1}^{n}{aC_{n}H{n}}}{\sqrt{d}}
    \end{equation}

    CN : (0ºC y 1 atm  $ \mid 1 \ \mathrm{mol} = 22.4 \ \mathrm{L}$).
    
    \textbf{Condiciones estándar: (25ºC y 1 atm).}

    \paragraph*{Dosado}

    \begin{equation}
        F = \frac{m_f}{m_a}
    \end{equation}

    \begin{equation}
        F_R = \frac{F}{F_{e}}
    \end{equation}

    Se tiene 
    \begin{align*}
        F_{R} < 1 &\rightarrow \text{Mezcla rica en aire (pobre en combustible)} \\
        F_{R} = 1 &\rightarrow \text{Mezcla estequiométrica} \\
        F_{R} > 1 &\rightarrow \text{Mezcla rica en combustible (pobre en aire)}
    \end{align*}

    \begin{equation}
        \lambda = \frac{1}{F_R}
    \end{equation}

    \begin{align*}
        \lambda < 1 &\rightarrow \text{Mezcla rica en combustible (pobre en aire)} \\
        \lambda = 1 &\rightarrow \text{Mezcla estequiométrica} \\
        \lambda > 1 &\rightarrow \text{Mezcla rica en aire (pobre en combustible)}
    \end{align*}

    \begin{equation}
        E = \frac{21}{21 - O_2} \cdot 100
    \end{equation}

    \begin{align*}
        E < 100 &\rightarrow \text{Mezcla rica en combustible (pobre en aire)} \\
        E = 100 &\rightarrow \text{Mezcla estequiométrica (aire estequiométrico)} \\
        E > 100 &\rightarrow \text{Mezcla rica en aire (pobre en combustible)}
    \end{align*}

    \paragraph*{Calderas de condensación}

    La condiciones para que una caldera condense es: 
    \[
        T_{\text{retorno}} < T_{\text{rocío}}
    \]

    \paragraph*{Ecuación de los Gases Ideales}
    \begin{equation}
        PV = nRT
    \end{equation}

    de donde $R = 8.314 \ \mathrm{J/(mol \cdot K)}$.

    \paragraph*{IMPORTANTE}

    \begin{align}
        \text{Poder comburívoro} &= \frac{\text{Nm}^{3} \text{ de aire}}{\text{KWh}} \vphantom{\frac{\text{Nm}^{3} \text{humos secos}}{\text{KWh}}} \\
        \text{Poder fumígeno} &= \frac{\text{Nm}^{3} \text{humos secos}}{\text{KWh}}
    \end{align}

    \textcolor{red}{\textbf{¡OJO!: $Nm^3$ es $m^3$ a condiciones normales (0ºC y 1 atm).}} 
    

\end{multicols*}


\section*{BLOQUE 2--VAPOR}

\begin{multicols*}{2}
    \paragraph*{Termodinámica Fundamental del Vapor}
    
    \paragraph*{Definición de Presiones}
    Relación entre la presión absoluta, la relativa (manométrica) y la atmosférica.
    \begin{equation}
        P_{abs} = P_{rel} + P_{atm}
    \end{equation}
    
    Conversión de unidades de presión a altura de columna de agua (m.c.a).
    \begin{equation}
        1 \text{ bar} = 100 \text{ kPa} = 10,2 \text{ m.c.a.}
    \end{equation}
    
    \paragraph*{Entalpía y Calidad del Vapor}
    Cálculo de la entalpía total del vapor saturado ($h_g$) como suma de la entalpía del líquido saturado ($h_f$) y la entalpía de vaporización ($h_{fg}$).
    \begin{equation}
        h_g = h_f + h_{fg}
    \end{equation}
    
    Cálculo de la entalpía para vapor húmedo en función de la fracción seca ($x$).
    \begin{equation}
        h = h_f + x \cdot h_{fg}
    \end{equation}
    
    \paragraph*{Dinámica de Fluidos y Propiedades de Transporte}
    
    \paragraph*{Viscosidad}
    Relación entre la viscosidad cinemática ($\vartheta$), dinámica ($\mu$) y la densidad ($\rho$).
    \begin{equation}
        \vartheta = \frac{\mu}{\rho}
    \end{equation}
    
    \paragraph*{Ecuación de Continuidad y Velocidad}
    Cálculo de la velocidad del fluido ($V$) en función del caudal volumétrico ($Q$) y el diámetro interior ($D$).
    \begin{equation}
        V = \frac{4 \cdot Q}{\pi \cdot D^2}
    \end{equation}
    
    
    \paragraph*{Número de Reynolds ($\mathcal{R}$)}
    Criterio adimensional para determinar el régimen de flujo (Laminar si $\mathcal{R} < 2300$, Turbulento si $\mathcal{R} > 2300$).
    
    En función de la velocidad ($V$) y diámetro ($D$):
    \begin{equation}
        \mathcal{R} = 10^3 \frac{V \cdot D}{\vartheta}
    \end{equation}
    
    En función del caudal ($Q$):
    \begin{equation}
        \mathcal{R} = \frac{4 \cdot 10^6 \cdot Q}{3,6 \cdot \pi \cdot D \cdot \vartheta}
    \end{equation}
    
    \paragraph*{Pérdida de Carga e Hidráulica de Tuberías}
    
    \paragraph*{Carga Estática y Conversión}
    Relación de Bernoulli simplificada para la altura de carga ($h_f$) en función de la presión ($p$) y el peso específico ($\omega$) en $kg/m^3$.
    \begin{equation}
        h_f = \frac{p}{\omega}
    \end{equation}
    
    Conversión práctica para fluidos reales (donde $p$ está en $kg/cm^2$ y $h_f$ en metros).
    \begin{equation}
        h_f = 10000 \frac{p}{\omega}
    \end{equation}

    \paragraph*{Ecuacion de Prony}

    Es una ecuación empírica.

    \begin{equation}
        h_f=\frac{L}{D}\cdot \left(aV + bV^2\right)
    \end{equation}
    
    Donde $L/D$ es la relación entre la longitud y el diámetro de la tubería, $V$ es la velocidad del fluido por la tubería, y $a$ y $b$ son coeficientes empíricos que dependen del tipo de fluido y las condiciones de flujo.
    \paragraph*{Ecuación de Darcy-Weisbach}
    Cálculo de la pérdida de carga ($h_f$) en tuberías rectas.
    \begin{equation}
        h_f = f \cdot \frac{L}{D} \cdot \frac{V^2}{2g}
    \end{equation}
    Donde $f$ es el factor de fricción de Darcy.
    
    Expresión alternativa en función del caudal ($Q$):
    \begin{equation}
        h_f = f \cdot \frac{L}{D^5} \cdot \frac{8 \cdot Q^2}{\pi^2 g}
    \end{equation}
    
    \paragraph*{Ecuación de White-Colebrook}
    Ecuación implícita para determinar el factor de fricción ($f$) en régimen turbulento y de transición.
    \begin{equation}
        \frac{1}{\sqrt{f}} = -2 \log_{10} \left( \frac{\varepsilon/D}{3,7} + \frac{2,51}{\mathcal{R}\sqrt{f}} \right)
    \end{equation}

    Donde $\epsilon/D$ es la rugosidad relativa y $\epsilon$ es la rugosidad absoluta de la tubería.

    Caso particular: \textbf{Tuberías lisas}
    \begin{equation}
        \frac{1}{\sqrt{f}}= 2 \cdot log_{10} \left[ \mathcal{R} \cdot \sqrt{f} \right] - 0.8
    \end{equation}

    \paragraph*{Ecuación de Hazen-Williams}
        Para determinar la velocidad de agua en tubería circulares llenas, o conductos cerrados.
        En función del radio hidraúlico
        \begin{equation}
            V = 0.8494 \cdot C \cdot \left(\frac{r_h}{4}\right)^{0.63} \cdot S^{0.54}
        \end{equation}

        En función del diámetro
        \begin{equation}
            V = 0.2785 \cdot C \cdot \left(\frac{D}{4}\right)^{2.63} \cdot S^{0.54}
        \end{equation}

        \textit{Nota:} Esto es para instalaciones antiguas, esto yo no lo voy a utilizar.
    \paragraph*{Cálculo del Diámetro Mínimo (Sistema Iterativo)}
    Combinación de Darcy y Colebrook para iterar el diámetro ($D$) basado en una pérdida de carga admisible ($h_f$).

    \begin{equation}
        \frac{1}{\sqrt{f}}= -2 \log_{10} \left[ \frac{\varepsilon/D}{3,7} + \frac{2,51}{\mathcal{R}\sqrt{f}} \right]
    \end{equation}
    
    sustituyendo $\mathcal{R}$ y $f$ de Darcy:
    
    \begin{equation}
        \frac{1}{\sqrt{K D^5}} = -2 \log_{10} \left[ \frac{\varepsilon/D}{3,7} + \frac{2,51}{\frac{4Q}{\pi D \vartheta} \sqrt{K D^5}} \right]
    \end{equation}
    Donde el factor $K$ se define como:
    \begin{equation}
        K = h_f \cdot \frac{\pi^2 g}{8 \cdot L \cdot Q^2}
    \end{equation}
    
    \paragraph*{Cálculo de Cargas Térmicas y Puesta en Marcha}
    
    \paragraph*{Carga de Puesta en Marcha}
    Cantidad de vapor ($Q_v$) necesaria para calentar la masa metálica de la instalación.
    \begin{equation}
        Q_v = \frac{P \cdot (T_v - T_a) \cdot C_p}{C_L}
    \end{equation}
    Donde: $P$: Peso total de tubería y accesorios; $T_v, T_a$: Temperatura del vapor y ambiente; $C_p$: Calor específico del material (acero); $C_L$: Calor latente del vapor.
    
    Caudal horario de puesta en marcha ($Q$) en función del tiempo requerido ($t_r$ en minutos).
    \begin{equation}
        Q = Q_v \cdot \frac{60}{t_r}
    \end{equation}
    
    \section*{Recuperación de Condensados y Vapor Flash}
    
    \paragraph*{Revaporizado (Flash Steam)}
    Porcentaje de vapor flash generado por caída de presión, basado en el balance de entalpías.
    \begin{equation}
        \%\text{vapor flash} = \frac{h_{lp1} - h_{lp2}}{h_{vp2}-h_{lp2}} \cdot 100
    \end{equation}
    Donde: $h_{lp1}$: entalpía específica del líquido a la presión de la red de vapor; $h_{lp2}$: entalpía específica del líquido a la presión de la red de condensados; $h_{vp2}$: entalpía específica del vapor saturado a la presión de la red condensados.
    
    \section*{Transferencia de Calor y Aislamiento}
    
    \paragraph*{Flujo de Calor en Tuberías Cilíndricas}
    Densidad de flujo de calor ($q$) a través de una pared cilíndrica hueca.
    \begin{equation}
        q = \frac{T_{s,int} - T_{s,ext}}{R_j}
    \end{equation}

    Donde $T_{s,int}$ y $T_{s,ext}$ son las temperaturas en la superficie interna y externa de la tubería, respectivamente, y $R_j$ es la resistencia térmica linealde una capa cilíndrica (aislamiento) que se calcula como: 
    \begin{equation}
        R_j = \frac{\ln(D_e / D_i)}{2 \pi \lambda_j} \hspace{0.2cm} \text{[m$\cdot$ K/W]}
    \end{equation}
        
    Donde:$D_e$: Diámetro exterior de la capa [m]; $D_i$: Diámetro interior de la capa [m] y $\lambda_j$: Conductividad térmica del material [W/m·K].

    Pared multicapa

    \begin{equation}
        q = \frac{T_{s,int} - T_{s,ext}}{R'} \hspace{0.2cm} \text{[W/m]}
    \end{equation}

    donde $R'$ es la resistencia térmica total de la pared multicapa

    \begin{equation}
        R' = \frac{1}{2\pi} \sum_{j=1}^{n} {\left(\frac{1}{\lambda_j} \cdot \ln\frac{D_{e,j}}{D_{i,j}}\right)} \hspace{0.2cm} \text{[m$\cdot$ K/W]}
    \end{equation}


    \paragraph*{Coeficiente superficial de transmisión del calor}

        \begin{equation}
            h_s = h_r + h_{cv}
        \end{equation}

        Donde $h_r$ es la parte radiante y $h_{cv}$ la parte convectiva.

        La parte radiante viene dada por: 

        \begin{equation}
            h_r = a_r \cdot C_r \hspace{0.2cm} \text{[W/m$^2$K]}
        \end{equation}

        $C_r$ es el coeficiente de radiación en $W/m^2 \cdot K^4$ [producto $\epsilon$ y cte. Stephan Bolzman ($5.67 \cdot 10^{-8}$)] y $a_r$ es el factor de temperatura en $K^3$.

        siendo $a_r$ el factor de temperatura:

        \begin{equation}
            a_r = \frac{T_{1}^4 - T_{2}^4 }{T_1 - T_2} \hspace{0.2cm} \text{[$K^3$]}
        \end{equation}
    

    \paragraph{Parte convectiva del coeficiente superficial $h_{cv}$}

        En el interior del edificio encontramos:

        \begin{itemize}
            \item Tuberías VERTICALES:
                \begin{itemize}
                    \item Convención libre laminar ($D_{e}^3 \cdot \Delta T \leq 10 m^3 \cdot K$): 
                    \begin{equation}
                        h_{cv} = 1.32 \cdot \sqrt[4]{\frac{\Delta T}{D_e}} \hspace{0.2cm} \text{[W/m$^2$K]}
                    \end{equation}
                    \item Convención libre turbulenta ($D_{e}^3 \cdot \Delta T \geq 10 m^3 \cdot K$): 
                    \begin{equation}
                        h_{cv} = 1.74 \cdot \sqrt[3]{\Delta T} \hspace{0.2cm} \text{[W/m$^2$K]}
                    \end{equation}
                \end{itemize}

            \item Tuberías HORIZONTALES: 
                \begin{itemize}
                    \item Convención libre laminar ($D_{e}^3 \cdot \Delta T \leq 10 \hspace {0.1cm} m^3 \cdot K$): 
                    \begin{equation}
                        h_{cv} = 1.25 \cdot \sqrt[4]{\frac{\Delta T}{D_e}} \hspace{0.2cm} \text{[W/m$^2$K]}
                    \end{equation}
                    \item Convención libre turbulenta ($D_{e}^3 \cdot \Delta T \geq 10 \hspace{0.1cm} m^3 \cdot K$): 
                    \begin{equation}
                        h_{cv} = 1.21 \cdot \sqrt[3]{\Delta T} \hspace{0.2cm} \text{[W/m$^2$K]}
                    \end{equation}
                \end{itemize}
        \end{itemize}

        En el exterior del edificio tanto para tuberías HORIZONTALES como VERTICALES:

            \begin{itemize}
                \item Convención libre laminar ($ v \cdot D_{e}^3 \leq 8.55 \cdot 10^{-3} \hspace {0.1cm} m^2/s \cdot K$):
                \begin{equation}
                    h_{cv} = \frac{8.1 \cdot 10^{-3}}{D_e} + 3.14 \cdot \sqrt[4]{\frac{v}{D_e}} \hspace{0.2cm} \text{[W/m$^2$K]}
                \end{equation}

                \item Convención libre turbulenta ($ v \cdot D_{e}^3 \geq 8.55 \cdot 10^{-3} \hspace {0.1cm} m^2/s \cdot K$):
                \begin{equation}
                    h_{cv} = 8.9 \cdot \frac{v^{0.9}}{D^{0.1}_e} \hspace{0.2cm} \text{[W/m$^2$K]}
                \end{equation}
            \end{itemize}
        
        \paragraph*{Resistencia Térmica Superficial}
            Definición de la resistencia superficial ($R_{se}$) como inversa del coeficiente de película exterior ($h_e$).
            \begin{equation}
                R_{se} = \frac{1}{h_e \cdot \pi \cdot D_e} \hspace{0.2cm} \text{[m$\cdot$ K/W]}
            \end{equation}
        
        \paragraph*{Transmitancia térmica}

            \begin{equation}
                U_l = \frac{q_l}{T_i - T_a} \hspace{0.2cm} \text{[W/m$\cdot$ K]}
            \end{equation}

            su inversa es la resistencia térmica lineal:
            \begin{equation}
                \frac{1}{U_l} = R = R_{lpi} + R_{lc} + R_{lpe}
            \end{equation}

            paredes cilíndricas (tuberías y aislamientos de tubería):

            \begin{equation}
                \frac{1}{U_l} = \frac{1}{h_i \pi D^2{i}} + R_{lc} + \frac{1}{h_e \pi D^2_e} \hspace{0.2cm} \text{[m$\cdot$ K/W]}
            \end{equation}

        \paragraph*{Temperatura superficial}

            \begin{equation}
                T_{se} = T_a + \frac{T_i - T_a}{\frac{h_e D_e}{2\lambda} \cdot \ln\frac{D_e}{D_i} + 1} \hspace{0.2cm} \text{[ºC]}
            \end{equation}

            \textit{Nota: Para una sola capa de aislante en paredes cilíndricas.}
            
    
    
\end{multicols*}

%%%%%%%%%%%%%%%%%%
%%%% BLOQUE 3 %%%%
%%%%%%%%%%%%%%%%%%


\section*{BLOQUE 3--GAS NATURAL}
\begin{multicols*}{2}

    \paragraph*{Presiones de Gas NATURAL}
    \begin{itemize}
        \item \textbf{Alta presión B (APB):} $> 16$ bar
        \item \textbf{Media presión B (MPB):} Entre 2 y 5 bar
        \item \textbf{Media presión A (MPA):} Entre 0.1 y 2 bar
        \item \textbf{Baja presión (BP):} $>0.1$ bar
    \end{itemize}

    \paragraph*{Características del Gas Natural}
    \begin{itemize}
        \item \textbf{Densidad relativa ($d_r$):} $0.62$
        \item \textbf{PCS:} $11 \text{ kWh/m}^3\text{(s)}$
    \end{itemize}

    \textcolor{red}{\textbf{¡ATENCIÓN!: Los valores anteriores están referidos a condiciones estándar (25 \unit{\celsius} y 1 atm).}}

    \paragraph*{Consumo de un aparato ($Q_n$)}
    \begin{equation}
        Q_n = \frac{GC}{PCS}
    \end{equation}
    \begin{itemize}
        \item $Q_n$: Caudal nominal del aparato en $\text{m}^3\text{(s)/h}$.
        \item $GC$: Gasto calorífico del aparato referido al PCS en kW.
        \item $PCS$: Poder calorífico superior del gas en $\text{kWh/m}^3\text{(s)}$.
    \end{itemize}

    \paragraph*{Caudal simultáneo en vivienda ($Q_{si}$)}
    \begin{equation}
        Q_{si} = A + B + \frac{C + D + \dots + N}{2}
    \end{equation}
    \begin{itemize}
        \item $Q_{si}$: Caudal máximo de simultaneidad en $\text{m}^3\text{(s)/h}$.
        \item $A, B$: Caudales de los dos aparatos de mayor consumo en $\text{m}^3\text{(s)/h}$.
        \item $C, D, \dots, N$: Caudales del resto de aparatos en $\text{m}^3\text{(s)/h}$.
    \end{itemize}

    \paragraph*{Caudal en instalaciones comunes ($Q_{sc}$)}
    \begin{equation}
        Q_{sc} = \sum Q_{si} \cdot S_n
    \end{equation}
    \begin{itemize}
        \item $Q_{sc}$: Caudal máximo de simultaneidad de la instalación común en $\text{m}^3\text{(s)/h}$.
        \item $Q_{si}$: Caudal máximo de simultaneidad de cada vivienda o local en $\text{m}^3\text{(s)/h}$.
        \item $S_n$: Factor de simultaneidad (función del número de viviendas).
    \end{itemize}

    \paragraph*{Potencia simultánea de la instalación ($P_{nsc}$)}
    \begin{equation}
        P_{nsc} = Q_{sc} \cdot PCS
    \end{equation}
    \begin{itemize}
        \item $P_{nsc}$: Potencia nominal de utilización simultánea en kW (o kcal/h).
        \item $Q_{sc}$: Caudal máximo de simultaneidad de la acometida interior o instalación común en $\text{m}^3\text{(s)/h}$.
        \item $PCS$: Poder calorífico superior en $\text{kWh/m}^3\text{(s)}$ (o $\text{kcal/m}^3\text{(s)}$).
    \end{itemize}

    \paragraph*{Longitud equivalente ($L_E$)}
    Para compensar pérdidas en accesorios (codos, válvulas, etc.):
    \begin{equation}
        L_E = L_{real} \cdot 1.20
    \end{equation}

    \paragraph*{Pérdida de presión (Fórmulas de Renouard)}
    Válidas siempre que la velocidad del gas no supere los $20 \text{ m/s}$.

    \textbf{Renouard Lineal ($P_{rel} \leq 100 \text{ mbar}$):}
    \begin{equation}
        \Delta P = 23200 \cdot d_r \cdot L_E \cdot Q^{1.82} \cdot D^{-4.82}
    \end{equation}
    \begin{itemize}
        \item $\Delta P$: Diferencia de presión entre inicio y final en mbar.
        \item $d_r$: Densidad relativa del gas.
        \item $L_E$: Longitud equivalente del tramo en m.
        \item $Q$: Caudal en $\text{m}^3\text{(s)/h}$.
        \item $D$: Diámetro interior de la conducción en mm.
    \end{itemize}

    \textbf{Renouard Cuadrática ($P_{rel} > 100 \text{ mbar}$):}
    \begin{equation}
        {P_1}^2 - {P_2}^2 = 48.6 \cdot d_r \cdot L_E \cdot Q^{1.82} \cdot D^{-4.82}
    \end{equation}
    \begin{itemize}
        \item $P_1, P_2$: Presiones absolutas (efectiva + atmosférica) al inicio y final en \textbf{bar(a)}.
    \end{itemize}

    Mejor trabajar con:
    \begin{equation}
        {P_{\text{inicial}}}^2 - {P_{\text{final}}}^2 = 48.6 \cdot d_r \cdot L_E \cdot Q^{1.82} \cdot D^{-4.82}
    \end{equation}

    \paragraph*{Velocidad máxima del gas ($v$)}
    \begin{equation}
        v = 354 \cdot Q \cdot P^{-1} \cdot D^{-2}
    \end{equation}
    \begin{itemize}
        \item $v$: Velocidad del gas en m/s.
        \item $Q$: Caudal en $\text{m}^3\text{(s)/h}$.
        \item $P$: Presión absoluta al final del tramo en \textbf{bar(a)}.
        \item $D$: Diámetro interior de la conducción en mm.
    \end{itemize}

\end{multicols*}

%%%%%%%%%%%%%%%%%%
%%%% BLOQUE 4 %%%%
%%%%%%%%%%%%%%%%%%

\section*{BLOQUE 4--AIRE COMPRIMIDO}

\begin{multicols*}{2}

    \paragraph*{Trabajo de compresión con compresión isotérmica}
    \begin{equation}
        W = P_1 \cdot V_1 \cdot \ln\left(\frac{P_2}{P_1}\right)
    \end{equation}

    \paragraph*{Trabajo de compresión con compresión isentrópica}

    \begin{equation}
        W = \frac{k}{k-1} \cdot P_1 \cdot V_1 \cdot \left( P_2 \cdot V_2 - P_1 \cdot V_1 \right)
    \end{equation}

    \begin{minipage}{\linewidth}
        \begin{align*}
            W &= \text{Trabajo de compresión} \ \mathrm{[J]} \\
            k &= \text{Exponente isentrópico} \ K \approx 1.3 - 1.4 \\
            P_1, V_1 &= \text{Presión y volumen iniciales} \\
            P_2, V_2 &= \text{Presión y volumen finales}
        \end{align*}
    \end{minipage}

    \paragraph*{Cantidad de aire de ventilación}

    \begin{equation}
        q_v = \frac{P_v}{0.92 \cdot \Delta T}
    \end{equation}

    \begin{minipage}{\linewidth}
        \begin{align*}
            q_v &= \text{Cantidad de aire de ventilación} \ \mathrm{[m^3/h]} \\
            P_v &= \text{flujo térmico} \ \mathrm{[W]} \\
            \Delta T &= \text{Aumento de temperatura permitido} \ \unit{\celsius}
        \end{align*}
        
    \end{minipage}

    \section*{Dimensionado de la red de aire comprimido}

    \begin{equation}
        h_f = f \cdot \frac{L}{D^5} \cdot \frac{8 \cdot Q^2}{\pi^2 g}
    \end{equation}

    \begin{minipage}{\linewidth}
        \begin{align*}
            h_f &= \text{Pérdida de carga debida a la} \\
            &\quad \text{fricción} \ \mathrm{[m.c.a.c]} \\
            f &= \text{Factor de fricción de Darcy} \\
            L &= \text{Longitud de la tubería} \ \mathrm{[m]} \\
            D &= \text{Diámetro de la tubería} \ \mathrm{[m]} \\
            Q &= \text{Caudal volumétrico} \ \mathrm{[m^3/s]} \\
            g &= 9.81 \ \mathrm{m/s^2}
        \end{align*}
    \end{minipage}

    \paragraph*{Ecuación de White-Colebrook}

    \begin{equation}
        \frac{1}{\sqrt{f}} = -2 \log_{10} \left( \frac{\varepsilon/D}{3,7} + \frac{2,51}{\mathcal{R}\sqrt{f}} \right)
    \end{equation}

    \paragraph{Simplificación(para tubos rectos)}

    \begin{equation}
        \Delta P = 450 \cdot \frac{Q_{c}^{1.85} \cdot L}{D^{5} \cdot P}
    \end{equation}

    \begin{minipage}{\linewidth}
        \begin{align*}
            \Delta P &= \text{Pérdida de presión} \ \mathrm{[bar]} \\
            Q_c &= \text{Caudal de aire} \ \mathrm{[L/s]} \\
            L &= \text{Longitud eq. de tubería recta} \ \mathrm{[m]} \\
            D &= \text{Diámetro interno de la tubería} \ \mathrm{[mm]} \\
            P &= \text{Presión absoluta en cabeza de} \\
            &\quad \text{distribución} \ \mathrm{[bar(a)]}
        \end{align*}
    \end{minipage}

    \paragraph*{Volúmen de un Acumulador}

    \begin{equation}
        V = \frac{0.25 \cdot q_c \cdot P_1 \cdot T_0}{f_{\max}\cdot \left(P_c - P_d\right) \cdot T_1}
    \end{equation}

    \begin{minipage}{\linewidth}
        \begin{align*}
            V &= \text{Volumen del acumulador} \ \mathrm{[m^3]} \\
            q_c &= \text{Caudal de aire} \ \mathrm{[L/s]} \\
            P_1 &= \text{Presión de aspiración} \ \mathrm{[bar(a)]} \\
            T_1 &= \text{Temp. aire en compresor} \ \mathrm{[K]} \\
            T_0 &= \text{Temp. aire en receptor} \ \mathrm{[K]} \\
            P_c &= \text{Presión llenado acumulador} \ \mathrm{[bar(a)]} \\
            P_d &= \text{Presión descarga acumulador} \ \mathrm{[bar]} \\
            P_c - P_d &= \text{Diferencia de presión entre} \\ 
            &\quad \text{carga y descarga} \ \mathrm{[bar]} \\
            f_{\max} &= \text{frecuencia máxima de carga} \ \mathrm{[s^{-1}]} 
        \end{align*}
    \end{minipage}

    Para el cálculo de receptores aplicar la siguiente relación:

    \begin{equation}
        V = \frac{q \cdot t}{P_1 - P_2} = \frac{L}{P_1 - P_2}
    \end{equation}

    \begin{minipage}{\linewidth}
        \begin{align*}
            V &= \text{Volumen de aire del receptor} \ \mathrm{[m^3]} \\
            q &= \text{Caudal de aire durante la fase} \\
            &\quad \text{de llenado} \ \mathrm{[m^3/s]} \\
            t &= \text{Tiempo de descarga} \ \mathrm{[s]} \\
            P_1 &= \text{Presión normal} \\
            &\quad \text{de trabajo en la red} \ \mathrm{[bar]} \\
            P_2 &= \text{mínima presión admisible en} \\
            &\quad \text{la red} \ \mathrm{[bar]} \\
            L &= q \cdot t = \text{requerimiento de aire de} \\
            &\quad \text{llenado} \ \mathrm{[l / \text{ciclo de trabajo}]}
        \end{align*}
    \end{minipage}

    \paragraph*{Eliminación del agua}

    \begin{align*}
        f_3 = f_1 - f_2 \\
        f_2 = \frac{1 \cdot h_{sac} \cdot Q}{100 \cdot f_1} \\
        f_1 = \frac{H_r \cdot h_{saa} \cdot Q}{100}
    \end{align*}

    \begin{minipage}
    {\linewidth}
        \begin{align*}
            f_1 &= \text{Cantidad de agua remanente en el } \\
            &\quad \text{aire después de condensar} \ \mathrm{[g/s]} \\
            f_2 &= \text{Cantidad de agua en el aire } \\
            &\quad \text{antes de condensar} \ \mathrm{[g/s]} \\
            f_3 &= \text{Cantidad de condensación de agua} \ \mathrm{[g/s]} \\
            H_r &= \text{Humedad relativa del aire aspirado}  \ \mathrm{[\%]} \\
            h_{saa} &= \text{Humedad de saturación del aire aspirado} \ \mathrm{[g/l]}\\
            h_{sac} &= \text{Humedad de saturación del aire comprimido} \\ 
            &\quad \text{a la Tª de condensación} \ \mathrm{[\%]}\\
            Q &= \text{Caudal de aire del compresor} \ \mathrm{[l/s]}
        \end{align*}
    \end{minipage}

\end{multicols*}

\section*{BLOQUE 5--GLPs}
\begin{multicols*}{2}

    \paragraph*{Índice de Wobbe}
    \begin{equation*}
        W = \frac{P.C.S.}{\sqrt{d}}
    \end{equation*}

    existe otra expresión con el poder calorífico superior
    \begin{equation*}
        H = \frac{H_s}{\sqrt{d}}
    \end{equation*}

    En la práctica 
    \begin{equation*}
        W' = W \times K_1 \times K_2
    \end{equation*}

    \begin{minipage}{\linewidth}
        \begin{align*}
            W' &= \text{Índice de Wobbe corregido} \\
            W &= \text{Índice de Wobbe} \\
            K_1 , K_2 &= \text{Factores}
        \end{align*}
    \end{minipage}

    $K_1$ gases primera familia 
    \begin{equation*}
        H_2 - C_nH_m -2CO_2
    \end{equation*}

    $K_2$
    \begin{align*}
        1000& \frac{O_2}{P} \rightarrow \text{Primera familia} \\
        1000& \frac{CO +4O_2 - 0.5}{P} CO_2 \rightarrow \text{Segunda familia}
    \end{align*}

    \paragraph{Potencial de combustión}

    Primera familia 
    \begin{equation*}
        C = u \cdot \frac{H_2 + 0.3 CH_4 + 0.7CO +v \cdot \sum a \cdot C_nH_m}{\sqrt{d}}
    \end{equation*}

    Segunda familia 
    \begin{equation*}
        C = u \cdot \frac{H_2 + 0.7 CO_2 + 0.3 CH_4 +v \cdot \sum a \cdot C_nH_m}{\sqrt{d}}
    \end{equation*}

    donde 
    \begin{minipage}{\linewidth}
        \begin{align*}
            u &= \text{Coef. de corrección } (1000 \cdot O_2 / P) \\
            H_2, \dots &= \text{\% de componentes e} \\
            &\quad \text{hidrocarburos en el gas} \\
            v &= \text{Coeficiente según } H_2 \text{ o } W' \\
            &\quad \text{(1ª y 2ª fam. respectivamente)} \\
            C_n H_m &= \text{Hidrocarburos} \\
            d &= \text{Densidad relativa}
        \end{align*}
    \end{minipage}

    \paragraph{Autonomía}

    \begin{equation*}
        A = \frac{C}{c}
    \end{equation*}

    \begin{minipage}{\linewidth}
        \begin{align*}
            A &= \text{Autonomía} \\ 
            C &= \text{Capacidad de almacenamiento en kg} \\ 
            c &= \text{Consumo en kg/h} \\
        \end{align*}
    \end{minipage}

    \paragraph{Potencias}

    \begin{align*}
        P_s & = H_s \times Q \\
        P_i & = H_i \times Q \\
        P_u & = A \times \Delta T \\
        P_u & = P_s \times R_s \\
        & = H_s \times Q \times R_s \\
        & = P_i \times R_i \\
        & = H_i \times Q \times R_i 
    \end{align*}

    \begin{minipage}{\linewidth}
        \begin{align*}
            P_s &= \text{Potencia superior} \\
            P_i &= \text{Potencia inferior o gasto calorífico} \\
            P_u &= \text{Potencia útil} \\
            R_s &= \text{Rendimiento superior} \\
            R_i &= \text{Rendimiento inferior} \\
            Q &= \text{Caudal de gas} \\
            \Delta T &= \text{Incremento de temperatura ocasionado} \\
        \end{align*}
    \end{minipage}

    deducciones
    \begin{align*}
        \frac{H_i}{H_s} & = \frac{R_s}{R_i} \\
        \frac{P_i}{P_s} & = \frac{R_s}{R_i} \\
        \frac{R_s}{R_i} & = \frac{H_i}{H_s}
    \end{align*}

    \paragraph{Número de botellas}

    \begin{equation*}
        N = A \times \frac{C_d}{35}
    \end{equation*}

    \begin{minipage}{\linewidth}
        \begin{align*}
            N &= \text{Número de botellas} \\
            A &= \text{Autonomía} \\
            C_d &= \text{Capacidad de la botella en kg}
        \end{align*}
    \end{minipage}

    De aquí se deduce la autonomía como:
    \begin{equation*}
        A = \frac{\text{Gas almacenado}}{\text{Consumo diario}}
    \end{equation*}

    ``Cocinado''

    \begin{align*}
        N \text{ (bot.)} \times 35 \text{ (kg/bot.)} &= \text{Auton. (días)} \\
        &\quad \times \text{Cons. (kg/día)} 
    \end{align*}

    \paragraph{Caudal de simultaneidad}

    Siendo $Q_1 \geq Q_2 \geq Q_3 \geq \dots \geq Q_n$

    \begin{equation*}
        Q_{si} = Q_1 + Q_2 + \frac{1}{2} \cdot \left(Q_3 + Q_4 + \dots + Q_n\right)
    \end{equation*}

    \paragraph{Potencia útil aparato}

    \begin{equation*}
        P_u = \eta \cdot P_c
    \end{equation*}

    En los procesos de calentamiento:
    \begin{equation*}
        P_u = A \times Ce \times \Delta T
    \end{equation*}

    \paragraph{Rendimiento aparato}

    \begin{equation*}
        \eta = \frac{P_u}{P}
    \end{equation*}

    donde P es el gasto calorífico.

    \paragraph{Velocidad de un gas}

    \begin{equation*}
        v = \frac{Q}{S}
    \end{equation*}

    Para consistencia dimensional:

    \begin{equation*}
        v = 378.04 \times \frac{Q}{P \times D^2}
    \end{equation*}

    En baja presión (BP):
    \begin{equation*}
        v = 360 \times \frac{Q}{D^2}
    \end{equation*}

    \paragraph{Pérdidas de carga. Fórmulas de Renouard}
    
    Condición de aplicación:
    \begin{equation*}
        \frac{Q}{D} < 150 
    \end{equation*}

    Para media presión (MP) $0.05 \text{ bar} < P < 5 \text{ bar}$:

    \begin{equation*}
        P_A^2 - P_B^2 = 51.5 \cdot d_c \cdot L_c \cdot \frac{Q^{1.82}}{D^{4.82}}
    \end{equation*}

    \begin{minipage}{\linewidth}
        \begin{align*}
            P_A, P_B &= \text{Presiones abs. (origen y final) [bar a]} \\
            &\quad \text{(Relativa + 1.01325 bar)} \\
            d_c &= \text{Densidad corregida (Propano: 1.16; Butano: 1.44)} \\
            L_c &= \text{Longitud de cálculo de la conducción [m]} \\
            Q &= \text{Caudal de gas en el tramo [m}^3\text{/h]} \\
            D &= \text{Diámetro interior de la tubería [mm]}
        \end{align*}
    \end{minipage}

    Para baja presión $P < 0.05 bar$

    \begin{equation*}
        P_A - P_B = 25076 \cdot d_c \cdot L_c \cdot \frac{Q^{1.82}}{D^{4.82}}
    \end{equation*}

    \paragraph{Longitudes de cálculo}\mbox{}\\
    \textbf{BP}

    \begin{equation*}
        L_c = 1.05 \times L    
    \end{equation*}

    \textbf{MP}

    \begin{equation*}
        L_c = 1.05 \times L    
    \end{equation*}

    \paragraph{Pérdida de carga lineal}\mbox{}\\
    \textbf{BP}

    \begin{equation*}
        J = \frac{\left(P_A - P_B\right)}{L_c} = 25076 \times d_c \times \frac{Q^{1.82}}{D^{4.82}}
    \end{equation*}

    \textbf{MP}

    \begin{equation*}
        J = \frac{\left(P_A^2 - P_B^2\right)}{L_c} = 51.5 \times d_c \times \frac{Q^{1.82}}{D^{4.82}}
    \end{equation*}
        
    \paragraph{Pérdida de carga}

    \begin{equation*}
        J^c = \frac{\left(P_A ^{2} - P_B ^{2}\right)}{L_c}
    \end{equation*}

    \paragraph{Fórmulas de vaporización natural}

    \begin{equation*}
        Q = P \cdot S \cdot K \cdot \frac{T_e - T_g}{CLV}
    \end{equation*}

    \begin{minipage}{\linewidth}
        \begin{align*}
            Q &= \text{Caudal másico de vaporización [kg/h]} \\
            P &= \text{Coef. superficie mojada (20\%: 0.336; 30\%: 0.397)} \\
            S &= \text{Superficie del depósito [m}^2\text{]} \\
            K &= \text{Coef. transmisión de calor [kW/m}^2\text{·ºC]} \\
              &\quad \text{(Aéreos: 0.0116; Enterrados: 0.0086)} \\
            T_e &= \text{Temp. exterior mín. media prevista [ºC]} \\
            T_g &= \text{Temp. equilibrio líquido-gas del gas [ºC]} \\
            CLV &= \text{Calor latente de vaporización [0.11 kWh/kg]}
        \end{align*}
    \end{minipage}

    \paragraph{Carga}

    \begin{equation*}
        \text{Carga} = 0.85 \times V \times \rho = 0.85 \cdot 506 \cdot V = 430.1 \cdot V
    \end{equation*}

    Nota: La fórmula original indica $5.284 \cdot V$

    \paragraph{Presión máxima de servicio}

    \begin{equation*}
        P_mxs = \frac{PN}{Cs}
    \end{equation*}

    \paragraph{Potencia nominal de utilización simultánea}

    \begin{equation*}
        P_{\text{si}} = P_A + P_B + \frac{P_C + P_D + \dots + P_N}{3}
    \end{equation*}

    \paragraph{Caudal de utilización simultánea individual}

    \begin{equation*}
        Q_{\text{si}} = Q_A + Q_B + \frac{Q_C + Q_D + \dots + Q_N}{2}
    \end{equation*}

    \paragraph{Consumo diario}

    \begin{equation*}
        C_{\text{diario}} = Q \times T \text{ (m}^3\text{/día)}
    \end{equation*}

    \paragraph{Contenido mínimo del depósito}

    \begin{equation*}
        \text{Contenido} = C_{\text{diario}} \times N_{\text{días}}
    \end{equation*}

    \paragraph{Pérdida de carga linea por metro de conducción de cálculo}

    \begin{equation*}
        J = \frac{PCd}{Lc}
    \end{equation*}

    \paragraph{Espesor de la plataforma (Flotación)}
    \begin{equation*}
        e = \frac{1000 \cdot V - T}{2400 \cdot A \cdot L}
    \end{equation*}
    \begin{minipage}{\linewidth}
        \scriptsize
        \begin{align*}
            e &= \text{espesor de la plataforma (m)} \\
            V &= \text{volumen del depósito (m}^3\text{)} \\
            T &= \text{Tara (masa) del depósito (kg)} \\
            A &= \text{Ancho de la plataforma (m)} \\
            &\quad \text{(mínimo el diámetro del depósito)} \\
            L &= \text{Longitud de la plataforma (m)} \\
            &\quad \text{(mínimo la longitud del depósito)}
        \end{align*}
    \end{minipage}

    \paragraph{Espárragos de sujeción}

    \begin{equation*}
        T_u = \frac{E}{n}
    \end{equation*}

    \begin{minipage}{\linewidth}
        \scriptsize
        \begin{align*}
            T_u &= \text{Tensión unitaria} \\
            E &= \text{esfuerzo vertical ascendente} \\
            n &= \text{número de espárragos} \\
            g &= \text{aceleración de la gravedad}
        \end{align*}
    \end{minipage}

    Para un coeficiente de trabajo de $C_t = 4100 \cdot g$ del acero (tetracero, en kg/cm$^2$), resulta necesaria una sección de:
    \begin{equation*}
        S = \frac{T_u}{C_t}
    \end{equation*}

    \begin{equation*}
        D = \frac{\sqrt{1000 \times V -T}}{226.98}
    \end{equation*}

    \begin{minipage}{\linewidth}
        \scriptsize
        \begin{align*}
            D &= \text{diámetro del espárrago (mm)} \\
            V &= \text{volumen del depósito (m}^3\text{)} \\
            T &= \text{Tara del depósito (kg)} \\
        \end{align*}
    \end{minipage}

\end{multicols*}







\end{document}